% Chapter 1
\chapter{Introduction}
\label{Chapter1}
\lhead{Chapter 1. \emph{Introduction}}

\section{Background}
Geomagnetic observatories record the Earth's magnetic field at fixed sites over long periods. Most of the signal follows smooth daily and seasonal variation. Short-lived departures from that pattern may indicate local magnetic sources or disturbances that are of interest in geophysical monitoring and related applications. A practical system therefore needs three linked capabilities: reliable ingestion of time-stamped field data, a model of the expected quiet-time field, and a rule that separates normal variation from suspicious intervals.

Magnavis is a Python-based application developed in the laboratory of Prof.\ Saikat Ghosh at the Department of Physics, IIT Kanpur. It supports both interactive monitoring and offline analysis of observatory data. The current version focuses on the OBS2 sensor suite (three tri-axial channels treated as resultant magnitude for forecasting) and on comparing several forecasting and detection strategies under controlled benchmarks.

\section{Problem Statement}
Software that combines data visualisation, sequence forecasting, and anomaly detection in one workflow is still uncommon in open form. Many tools address only one of these steps, or they rely on fixed thresholds on raw field magnitude, which perform poorly when anomalies are extended in time or when the background itself drifts slowly. This thesis addresses the following questions:
\begin{enumerate}[label=(\alph*)]
    \item How can magnetic time series from a database or CSV export be processed in a repeatable pipeline from raw data to anomaly labels?
    \item Which forecasting models---statistical smoothers, attention bidirectional LSTM, pretrained GRU, and pretrained LSTM---best support residual-based detection under different ground-truth layouts?
    \item How sensitive is detection to the threshold multiplier $k$ in the rule $\text{threshold} = \mu + k\sigma$ applied to smoothed prediction residuals?
    \item When tri-axial components are available, how should the direction of the perturbation vector be computed and interpreted relative to a possible source?
\end{enumerate}

\section{Objectives}
The main objectives of this research are:
\begin{enumerate}[label=(\alph*)]
    \item To design and implement a modular desktop application that orchestrates data fetch, forecasting, anomaly detection, and visualisation for one or more sensors.
    \item To integrate recurrent predictors (GRU and LSTM) with optional pretraining on long historical CSV archives, and to document training and inference settings.
    \item To implement an anomaly detector based on the absolute difference between measured and forecast values, with an exponential moving average threshold controlled by $k$.
    \item To run a zero-historic, predict-only benchmark on short-duration, long-duration, and synthetic ground-truth scenarios and to report point-level recall, precision, and F1 across $k \in \{1,\ldots,5\}$.
    \item To explain the contrasting behaviour of GRU and LSTM on long-duration ground truth and to state practical deployment guidance.
    \item To describe perturbation-vector direction finding and triangulation where component data exist, together with known limitations.
\end{enumerate}

\section{Relevance}
Observatory-style monitoring benefits from methods that learn the diurnal baseline rather than assuming a constant level. The benchmark results show that the choice of predictor and threshold is not neutral: on long anomaly windows, simple smoothers applied directly to the raw field fail at moderate $k$, whereas neural predictors can maintain useful recall if the model family is chosen carefully. The software and evaluation protocol provided here can be reused for future sensors, retraining studies, and thesis extensions such as event-level metrics or direction calibration.

\section{Organisation of the Thesis}
The thesis is organised as follows:
\begin{enumerate}[label=(\alph*)]
    \item \textbf{Chapter 2} reviews magnetic anomaly detection, recurrent forecasting, and direction finding from tri-axial data.
    \item \textbf{Chapter 3} describes the Magnavis architecture and main modules.
    \item \textbf{Chapter 4} explains data acquisition and the residual-based anomaly detector.
    \item \textbf{Chapter 5} presents the GRU and LSTM predictors, the offline benchmark design, and the main experimental results.
    \item \textbf{Chapter 6} summarises the user interface and visualisation features.
    \item \textbf{Chapter 7} concludes and outlines future work.
\end{enumerate}
